%% Chapitre 02 — Installation et premier lancement
%% RÈGLE : uniquement des commandes sémantiques bookforge. Aucun \chapter,
%% \textbf, \begin{lstlisting}, \vspace... brut.

\chapitre{Installation et premier lancement}

\introChapitre{
  On passe à la pratique. Ce chapitre, volontairement daté car les commandes
  évoluent vite, t'amène d'un terminal vide à une première session Claude Code
  qui répond dans ton dépôt. On vérifie l'installation, l'authentification, et
  on lance un tout premier échange anodin pour prendre ses repères.
}

\paragraphe{Au chapitre précédent, on a posé le modèle mental : un agent
perçoit, agit, observe, et tu restes responsable du résultat. Bien. Maintenant
qu'on sait ce qu'est un agent, le plus simple pour le comprendre vraiment, c'est
d'en lancer un. Ce chapitre est le plus terre-à-terre du livre : pas de grande
idée, juste une suite d'étapes pour passer d'un terminal nu à un agent qui te
répond dans ton propre dépôt.}

\begin{avertissement}
Ce chapitre cite des commandes, des noms de paquets et des versions. C'est
exactement le genre de détail qui périme vite : un outil installé autrement
dans six mois, un drapeau renommé, une procédure d'authentification revue. Tiens
les \important{principes} (un gestionnaire de paquets, une vérification de
version, une authentification, un premier lancement) pour acquis ; vérifie les
\important{commandes exactes} dans la documentation officielle au moment où tu
lis. En cas de doute, \code{claude --help} et la page d'installation en ligne
sont les seules sources à jour.
\end{avertissement}

\section{Prérequis : terminal, compte Anthropic}

\paragraphe{Trois choses suffisent. Un terminal dans lequel tu es à l'aise,
parce que Claude Code vit en ligne de commande et pas dans une fenêtre à part.
Un compte Anthropic, une Console Claude ou un accès entreprise pris en charge,
qui servira à t'authentifier au premier lancement. Et, côté système, un
environnement de développement classique : Git, un shell et les outils habituels
de ton projet. Node.js n'est nécessaire que si \emphase{ton projet} ou un
serveur MCP que tu installes en dépend, pas pour installer Claude Code dans le
cas général.}

\paragraphe{Côté système, macOS, Linux, WSL et Windows natif sont pris en charge
par les méthodes d'installation officielles. Sur Windows, installe aussi Git for
Windows si tu veux que l'outil dispose d'un environnement Bash familier ; sinon,
il peut s'appuyer sur PowerShell. Tu sais déjà coder, donc je ne t'explique ni
ce qu'est un terminal ni comment installer Git : assure-toi simplement d'avoir
un environnement maintenu.}

\begin{extraitcode}{bash}
# Vérifier les outils de base avant de commencer
git --version
\end{extraitcode}

\paragraphe{Si cette commande renvoie un numéro de version, le terrain est prêt.
Si ton projet utilise Node, Python, Go ou un autre runtime, vérifie-le aussi :
Claude Code s'appuie sur les commandes de ton dépôt pour tester, construire et
formater. L'agent n'a pas besoin d'un environnement magique ; il a besoin du
même environnement que toi.}

\begin{note}
La distribution exacte de Claude Code peut évoluer. Le réflexe ne change pas :
ouvre la page d'installation officielle, suis la méthode recommandée pour ton
système, puis reviens ici pour la suite.
\end{note}

\section{Installer Claude Code et vérifier la version}

\paragraphe{L'installation tient en général en une commande. À l'état 2026, les
méthodes recommandées sont l'installateur natif, Homebrew sur macOS, et WinGet
sur Windows. Choisis celle qui correspond à ton poste ; les commandes suivantes
sont des repères, pas une promesse éternelle.}

\begin{extraitcode}{bash}
# macOS, Linux ou WSL : installateur natif
curl -fsSL https://claude.ai/install.sh | bash

# macOS avec Homebrew
brew install --cask claude-code

# Windows avec WinGet
winget install Anthropic.ClaudeCode

# Vérifier que la commande est disponible et connaître sa version
claude --version
\end{extraitcode}

\paragraphe{Cette dernière ligne n'est pas une formalité. La version que tu vois
détermine quelles fonctionnalités existent et comment elles se comportent. Quand
tu suivras un tutoriel ou que tu liras un message d'aide, c'est ce numéro qui te
dit si ce que tu lis s'applique à toi. Prends l'habitude de le vérifier, surtout
après une mise à jour.}

\begin{astuce}
Garde Claude Code à jour, mais pas en aveugle. Avant une montée de version sur
un projet sérieux, jette un œil aux notes de version : une commande peut avoir
changé de nom ou de comportement. L'installateur natif se met à jour
automatiquement ; avec Homebrew ou WinGet, utilise le mécanisme de mise à jour
du gestionnaire choisi. Dans tous les cas, \code{claude --version} confirme le
résultat.
\end{astuce}

\section{Authentification et configuration minimale}

\paragraphe{Installé ne veut pas dire connecté. Au premier lancement, Claude
Code a besoin de savoir qui tu es, pour rattacher ton usage à ton compte
Anthropic. Tu lances la commande, et l'outil t'accompagne : il ouvre en général
ton navigateur pour que tu te connectes, puis récupère un jeton qu'il stocke
localement.}

\begin{extraitcode}{bash}
# Premier lancement : déclenche l'authentification si nécessaire
claude
\end{extraitcode}

\paragraphe{Tu n'as quasiment rien à configurer pour démarrer. C'est volontaire :
les réglages fins (permissions, mémoire projet, outils externes) viendront plus
tard, quand tu en auras besoin. Pour l'instant, l'objectif est minimal : que
l'agent réponde. La configuration, elle, se construit au fil des chapitres
suivants, pas d'un coup ici.}

\begin{avertissement}
Le détail de l'authentification est précisément ce qui change le plus souvent :
connexion par navigateur, clé d'API à coller, choix entre un abonnement et une
facturation à l'usage. Ne mémorise pas une procédure figée. Retiens le principe :
au premier lancement, tu prouves ton identité une fois, et l'outil garde un jeton
pour les fois suivantes. Si l'écran que tu vois ne ressemble pas à ce qui est
décrit ici, fie-toi à ce que l'outil t'affiche, pas à ce livre.
\end{avertissement}

\section{Lancer une première session dans un dépôt jouet}

\paragraphe{Ne lance jamais ton tout premier agent sur ton projet le plus
précieux. Crée un \terme{dépôt jouet} : un petit dépôt sans enjeu, dans lequel
tu peux tout casser sans conséquence. C'est là qu'on ouvre notre première
\terme{session} : une conversation continue avec l'agent, avec son propre
contexte qui s'accumule au fil des échanges. Démarrer, reprendre ou
réinitialiser une session, c'est piloter ce que l'agent se rappelle.}

\begin{extraitcode}{bash}
# Créer un petit dépôt sans enjeu, puis y lancer Claude Code
mkdir bac-a-sable && cd bac-a-sable
git init
echo "# Bac à sable" > README.md

# Démarrer une session dans CE répertoire
claude
\end{extraitcode}

\paragraphe{Un point qui surprend souvent les débutants : l'agent travaille dans
le répertoire d'où tu l'as lancé. Son contexte de départ, c'est ce dossier-là.
Lancer Claude Code depuis le bon endroit n'est pas un détail, c'est ce qui
définit le périmètre de ce qu'il voit et de ce qu'il peut toucher.}

\paragraphe{Pour ce premier contact, demande quelque chose d'anodin et de
vérifiable. Pas « refais mon projet », mais une question dont tu connais déjà la
réponse, juste pour voir l'agent percevoir puis répondre.}

\begin{extraitcode}{text}
> Quels fichiers y a-t-il dans ce dépôt ? Résume-moi le README.
\end{extraitcode}

\paragraphe{Regarde ce qui se passe : l'agent ne devine pas, il va lire. Il
liste le répertoire, ouvre le README, et te répond à partir de ce qu'il a
réellement vu. C'est la boucle perception-action du chapitre précédent, sous tes
yeux, sur ton propre dossier. Rien de magique, et c'est tant mieux : tu peux
vérifier chaque affirmation.}

\section{Anatomie de l'interface : prompt, réponses, indicateurs}

\paragraphe{L'interface est sobre, et c'est une bonne nouvelle : il y a peu de
choses à apprendre. Trois zones reviennent en permanence.}

\begin{liste}
  \element{\important{Le prompt} : la ligne où tu tapes ta demande. C'est ton
  point d'entrée dans la session, l'équivalent de la barre d'adresse pour un
  navigateur.}
  \element{\important{Les réponses et les actions} : l'agent ne se contente pas
  d'écrire du texte. Il t'annonce les actions qu'il entreprend (lire un fichier,
  lancer une commande) et t'en montre le résultat. C'est là que tu suis ce qu'il
  fait, pas seulement ce qu'il dit.}
  \element{\important{Les indicateurs} : des repères discrets te signalent
  l'état courant — l'agent réfléchit, attend ton approbation, ou a terminé.
  Apprends à les lire : ils te disent à qui est la main, toi ou lui.}
\end{liste}

\paragraphe{Tu remarqueras vite que l'agent te demande parfois la permission
avant d'agir, surtout pour les actions qui modifient quelque chose. Garde ce
réflexe en tête sans t'y attarder maintenant : on consacrera un chapitre entier
à ces garde-fous. Pour l'instant, il suffit de savoir qu'un agent bien réglé
te demande avant de toucher à ce qui compte.}

\begin{astuce}
Le moyen le plus rapide de découvrir l'interface, ce n'est pas de tout lire,
c'est d'essayer. Tape une commande d'aide dès ta première session pour voir les
raccourcis et les commandes internes disponibles dans ta version.
\end{astuce}

\begin{extraitcode}{text}
> /help
\end{extraitcode}

\section{Quitter proprement et reprendre une session}

\paragraphe{Une session a un début et une fin. Tu la quittes proprement avec une
commande de sortie ou le raccourci d'interruption habituel de ton terminal. Rien
ne se perd côté code : l'agent travaille sur tes vrais fichiers, pas dans un
bac fermé qui s'évapore.}

\begin{extraitcode}{text}
> /exit
\end{extraitcode}

\paragraphe{Ce qui disparaît en revanche, c'est le contexte accumulé pendant
l'échange. Une nouvelle session repart sans mémoire de la précédente. C'est
souvent ce que tu veux — un contexte propre évite que de vieux échanges
parasitent une nouvelle tâche — mais pas toujours. Quand tu as construit un
contexte utile et que tu veux le retrouver, tu peux \important{reprendre} la
session précédente plutôt que d'en ouvrir une vierge.}

\begin{extraitcode}{bash}
# Reprendre la session précédente plutôt que repartir de zéro
claude --continue
\end{extraitcode}

\begin{note}
Le nom exact de l'option pour reprendre une session, comme tout le reste de ce
chapitre, peut changer. Le principe est durable : une session porte un contexte,
le quitter le met de côté, le reprendre le ramène. Savoir choisir entre
« repartir propre » et « reprendre » est une vraie compétence ; on y reviendra
quand le contexte deviendra une ressource à gérer.
\end{note}

\resumeChapitre{
  \element{\important{Prérequis} : un terminal, Git et un compte Anthropic ou Claude suffisent pour démarrer — vérifie ton environnement de projet avant d'installer.}
  \element{\important{Installation et version} : installe Claude Code avec la méthode officielle de ton système et contrôle immédiatement la version affichée ; c'est elle qui détermine les fonctionnalités disponibles.}
  \element{\important{Dépôt jouet} : lance ta première session sur un projet sans enjeu, jamais sur ton code de production, pour apprivoiser l'outil sans risque.}
  \element{\important{Répertoire de lancement} : l'agent travaille dans le dossier depuis lequel tu le lances — lancer depuis le bon endroit définit son périmètre d'action.}
  \element{\important{Session} : quitter la session efface le contexte accumulé ; reprends-la avec \code{--continue} quand tu veux retrouver l'historique précédent.}
}

\paragraphe{Voilà : tu as installé Claude Code, ouvert une session, vu l'agent
lire ton dépôt et répondre, puis quitté proprement. C'était volontairement
inoffensif — un README résumé ne casse rien. Au chapitre suivant, on arrête les
échauffements : on confie à l'agent une vraie petite tâche sur du vrai code, et
on apprend à relire ce qu'il propose avant de l'accepter.}
